\section*{第 \ref{ch:b3:radicals-carbenes} 章——自由基、卡宾与重排}

\begin{solution}{exo:b3:radicals-carbenes:1}
丙烷有六个伯氢和两个仲氢。氯代：$6 \times 1 : 2 \times 3.9 = 6 : 7.8$，
即 43\,\% 的 1-氯丙烷和 57\,\% 的 2-氯丙烷。溴代：$6 : 164$，4\,\% 和
96\,\%。
\end{solution}

\begin{solution}{exo:b3:radicals-carbenes:2}
$\ce{CH2}$：三重态；$\ce{CCl2}$：单重态（氯的孤对电子抬高了空的 p 轨道）。
二氯卡宾立体专一地加成，得到 \textit{trans}-1,1-二氯-2,3-二甲基环丙烷。
\end{solution}

\begin{solution}{exo:b3:radicals-carbenes:3}
\textit{cis}-1,2-二乙基环丙烷：类卡宾在一步中从某一面加成。
\end{solution}

\begin{solution}{exo:b3:radicals-carbenes:4}
环己酮：$\varepsilon$-己内酯（氧杂环庚-2-酮）。苯乙酮：乙酸苯酯，
$\ce{CH3CO2C6H5}$；苯基优先于甲基迁移。
\end{solution}

\begin{solution}{exo:b3:radicals-carbenes:5}
九个伯氢和一个叔氢。氯代：$9 : 5.1$，64\,\% 的 1-氯-2-甲基丙烷
和 36\,\% 的 2-氯-2-甲基丙烷。溴代：$9 : 1600$，0.6\,\% 和 99.4\,\%。
\end{solution}

\begin{solution}{exo:b3:radicals-carbenes:6}
$k_c = 3.0 \times k_H[\mathrm{Bu_3SnH}] = 3.0 \times \num{2.4e6} \times 0.020 = \qty{1.4e5}{s^{-1}}$.
\end{solution}

\begin{solution}{exo:b3:radicals-carbenes:7}
对 C5 进攻：5-exo-trig，有利；对 C6 进攻：6-endo-trig，按规则有利，但
此处较慢。4-exo-tet：有利。5-endo-trig：不利。5-exo-dig：有利。
3-exo-dig：不利。5-endo-dig：有利。
\end{solution}

\begin{solution}{exo:b3:radicals-carbenes:8}
失去水后生成叔苄基正离子；相邻碳上的苯基
和甲基均可迁移，迁移能力较强的苯基发生迁移。
氧稳定化的正离子失去质子，得到 3,3-二苯基丁-2-酮。
\end{solution}

\begin{solution}{exo:b3:radicals-carbenes:9}
库尔提斯重排：$\ce{C6H5CON3 -> C6H5NCO + N2}$。与水反应时，氨基甲酸失去 $\ce{CO2}$，生成苯胺
（苯胺与更多异氰酸酯相遇时会生成一些二苯脲）。与乙醇反应则得到
\textit{N}-苯基氨基甲酸乙酯，$\ce{C6H5NHCO2C2H5}$。
\end{solution}

\begin{solution}{exo:b3:radicals-carbenes:10}
氯：伯氢 $421.7 - 432 = \qty{-10}{kJ/mol}$，叔氢 $405.4 - 432 = \qty{-27}{kJ/mol}$；溴：$+56$ 和
$+39$。两步氯抽氢均为放热，过渡态出现较早，几乎感受不到
\qty{16}{kJ/mol} 的差别；两步溴抽氢均为吸热，过渡态出现较晚，其活化能差占其中大部分：
其活化能差占其中大部分：
$\eu^{16000/(8.314 \times 298)} \approx 600$，与观测比值的数量级相当。
\end{solution}

\begin{solution}{exo:b3:radicals-carbenes:11}
\textit{E} 肟（OH 与带甲基的 C2 呈反式）：C2 迁移，氮原子插入
其旁侧，得 7-甲基氮杂环庚-2-酮。\textit{Z} 肟：C6 迁移，得 3-甲基氮杂环庚-2-酮。在
贝克曼重排中由肟的几何构型决定，在拜耳–维利格
氧化中则由迁移能力决定，仲碳迁移：
7-甲基氧杂环庚-2-酮。
\end{solution}

\begin{solution}{exo:b3:radicals-carbenes:12}
环化分数 $k_c/(k_c + k_H c)$：0.91、0.49 和 0.09。对于 95\,\%：$k_Hc = k_c(1/0.95 - 1)$，于是
$c = 0.0526 \times \num{2.3e5}/\num{2.4e6} = \qty{5.0e-3}{mol/L}$。该浓度通过注射泵
缓慢加入氢化物来维持，或用催化量锡卤化物与还原剂，
使氢化物得以再生。
\end{solution}

\begin{solution}{pb:b3:radicals-carbenes:1}
\textbf{1.} $\ce{C6H10O + NH2OH -> C6H11NO + H2O}$。

\textbf{2.} 需要酸使加成中间体的羟基以水的形式脱去，
但酸过量则会使羟胺质子化，从而不再参与加成。

\textbf{3.} 否：与 C=N 碳相连的两个碳，
由于环的对称性而等价。

\textbf{4.} 在 \textit{E} 肟中，OH 指向远离带甲基的碳 C2 的一侧（与之呈反式
），在 \textit{Z} 肟中则指向它。

\textbf{5.} 连有电负性氧原子的肟氮，其翻转具有较高的
势垒。

\textbf{6.} 它使羟基质子化（或磺酸化），使其成为良好的
离去基团，即水分子。

\textbf{7.} 与离去基团呈反式的环 C--C 键由碳迁移到氮上，
氮由此插入环中，环原子数由六个变为七个。

\textbf{8.} 环状的腈正离子。

\textbf{9.} $\ce{C6H11NO -> C6H11NO}$：异构化，由肟转化为内酰胺。

\textbf{10.} 己内酰胺（氮杂环庚-2-酮），即 6-氨基己酸的七元
环状酰胺；其开环聚合得到尼龙-6。

\textbf{11.} \qty{98.0}{g/mol} 和 \qty{113.0}{g/mol}。

\textbf{12.} $\qty{1.0e6}{g}/\qty{98.0}{g/mol} = \qty{10204}{mol}$；$\times 0.98 \times 0.98 = \qty{9800}{mol}$；$\times \qty{113.0}{g/mol} = \qty{1.11e6}{g}$，即 \qty{1.11}{t}。

\textbf{13.} 肟：$\num{10204} \times 0.98 = \qty{1.00e4}{mol}$；$\ce{H2SO4}$：$\qty{1.30e4}{mol}$，可生成等摩尔数的
$\ce{(NH4)2SO4}$：$\qty{1.30e4}{mol} \times \qty{132.1}{g/mol} = \qty{1.7}{t}$，多于己内酰胺的质量。

\textbf{14.} $\ce{C6H10O + RCO3H -> C6H10O2 + RCOOH}$。

\textbf{15.} 过氧酸在羰基碳上形成的四面体加合物（即
克里吉中间体），一种过氧半缩酮。

\textbf{16.} 环上的 $\ce{CH2}$ 基团（两者等价）迁移到较近的那个过氧氧原子上；
羧酸根 $\ce{RCOO-}$ 离去。

\textbf{17.} $\varepsilon$-己内酯，即氧杂环庚-2-酮；其开环聚合得到
聚（$\varepsilon$-己内酯），一种可生物降解的聚酯。

\textbf{18.} 其唯一副产物为水，而过氧酸每生成一摩尔内酯便留下一摩尔
羧酸，且过氧酸本身储存时有危险。

\textbf{19.} 仲碳 C2 迁移：得 7-甲基氧杂环庚-2-酮。

\textbf{20.} 是：迁移的碳保持其构型。

\textbf{21.} 7-甲基氮杂环庚-2-酮（C2 呈反式并发生迁移）。

\textbf{22.} 3-甲基氮杂环庚-2-酮（C6 迁移）。

\textbf{23.} 贝克曼重排：由肟的构型决定（反式基团迁移）。
拜耳–维利格氧化：由迁移能力决定（取代较多的碳迁移）。

\textbf{24.} 其两个 $\alpha$ 碳等价：肟、内酰胺、内酯各一种。

\textbf{25.} 一吨环己酮约可得 \qty{1.11}{t} 己内酰胺。
\end{solution}
